% This must be in the first 5 lines to tell arXiv to use pdfLaTeX, which is strongly recommended.
\pdfoutput=1
% In particular, the hyperref package requires pdfLaTeX in order to break URLs across lines.

\documentclass[11pt]{article}

% Change "review" to "final" to generate the final (sometimes called camera-ready) version.
% Change to "preprint" to generate a non-anonymous version with page numbers.
\usepackage{acl}

% Standard package includes
\usepackage{times}
\usepackage{latexsym}
\usepackage{enumitem}
% For proper rendering and hyphenation of words containing Latin characters (including in bib files)
\usepackage[T1]{fontenc}
% For Vietnamese characters
% \usepackage[T5]{fontenc}
% See https://www.latex-project.org/help/documentation/encguide.pdf for other character sets

% This assumes your files are encoded as UTF8
\usepackage[utf8]{inputenc}

% This is not strictly necessary, and may be commented out,
% but it will improve the layout of the manuscript,
% and will typically save some space.
\usepackage{microtype}

% This is also not strictly necessary, and may be commented out.
% However, it will improve the aesthetics of text in
% the typewriter font.
\usepackage{inconsolata}

%Including images in your LaTeX document requires adding
%additional package(s)
\usepackage{graphicx}
\usepackage{multirow}
\usepackage{bm}
\usepackage{algorithm}
\usepackage{algpseudocode}
\usepackage{amsmath} % 用于公式支持
\usepackage{mathtools}
\usepackage{hyperref}
\usepackage{booktabs}
\usepackage{amssymb}
\usepackage{microtype}
% If the title and author information does not fit in the area allocated, uncomment the following
%
%\setlength\titlebox{<dim>}
%
% and set <dim> to something 5cm or larger.

\title{M\textsuperscript{3}HG: Multimodal, Multi-scale, and Multi-type Node Heterogeneous Graph for Emotion Cause Triplet Extraction in Conversations}

% Author information can be set in various styles:
% For several authors from the same institution:
% \author{Author 1 \and ... \and Author n \\
%         Address line \\ ... \\ Address line}
% if the names do not fit well on one line use
%         Author 1 \\ {\bf Author 2} \\ ... \\ {\bf Author n} \\
% For authors from different institutions:
% \author{Author 1 \\ Address line \\  ... \\ Address line
%         \And  ... \And
%         Author n \\ Address line \\ ... \\ Address line}
% To start a separate ``row'' of authors use \AND, as in
% \author{Author 1 \\ Address line \\  ... \\ Address line
%         \AND
%         Author 2 \\ Address line \\ ... \\ Address line \And
%         Author 3 \\ Address line \\ ... \\ Address line}



\author{
    Qiao~Liang\quad Ying~Shen\thanks{$\ $  Corresponding authors.}\quad Tiantian~Chen\quad Lin~Zhang \\
    Tongji University, Shanghai, China\\
    \texttt{\{2333091, yingshen, 2111287, cslinzhang\}@tongji.edu.cn}\thanks{\noindent 
This work was supported in part by the National Natural Science Foundation of China under Grant 62476202 and 62272343, in part by the Fundamental Research Funds for the Central Universities.}
}


%\author{
%  \textbf{First Author\textsuperscript{1}},
%  \textbf{Second Author\textsuperscript{1,2}},
%  \textbf{Third T. Author\textsuperscript{1}},
%  \textbf{Fourth Author\textsuperscript{1}},
%\\
%  \textbf{Fifth Author\textsuperscript{1,2}},
%  \textbf{Sixth Author\textsuperscript{1}},
%  \textbf{Seventh Author\textsuperscript{1}},
%  \textbf{Eighth Author \textsuperscript{1,2,3,4}},
%\\
%  \textbf{Ninth Author\textsuperscript{1}},
%  \textbf{Tenth Author\textsuperscript{1}},
%  \textbf{Eleventh E. Author\textsuperscript{1,2,3,4,5}},
%  \textbf{Twelfth Author\textsuperscript{1}},
%\\
%  \textbf{Thirteenth Author\textsuperscript{3}},
%  \textbf{Fourteenth F. Author\textsuperscript{2,4}},
%  \textbf{Fifteenth Author\textsuperscript{1}},
%  \textbf{Sixteenth Author\textsuperscript{1}},
%\\
%  \textbf{Seventeenth S. Author\textsuperscript{4,5}},
%  \textbf{Eighteenth Author\textsuperscript{3,4}},
%  \textbf{Nineteenth N. Author\textsuperscript{2,5}},
%  \textbf{Twentieth Author\textsuperscript{1}}
%\\
%\\
%  \textsuperscript{1}Affiliation 1,
%  \textsuperscript{2}Affiliation 2,
%  \textsuperscript{3}Affiliation 3,
%  \textsuperscript{4}Affiliation 4,
%  \textsuperscript{5}Affiliation 5
%\\
%  \small{
%    \textbf{Correspondence:} \href{mailto:email@domain}{email@domain}
%  }
%}


\begin{document}


\maketitle

\begin{abstract}
Emotion Cause Triplet Extraction in Multimodal Conversations (MECTEC) has recently gained significant attention in social media analysis, aiming to extract emotion utterances, cause utterances, and emotion categories simultaneously.
% However, there are many challenges in the MECTEC task.
% The first one is the scarcity of related datasets.
% Currently, there is only one published dataset for this task and its dialogue scenarios are highly uniform, which hampers the development of models in this field.
However, the scarcity of related datasets, with only one published dataset featuring highly uniform dialogue scenarios, hinders model development in this field.
% To address this issue, we construct a multimodal, multi-scenario Chinese dataset, namely \textbf{MECAD}, for the MECTEC task.
% The MECAD dataset comprises 989 conversations extracted from 56 different TV series and covers a broad range of social scenarios.
To address this, we introduce \textbf{MECAD}, the first multimodal, multi-scenario MECTEC dataset, comprising 989 conversations from 56 TV series spanning a wide range of dialogue contexts. 
% To the best of our knowledge, MECAD is the first Chinese multimodal dataset designed for the MECTEC task and related emotion cause analysis tasks. 
% In addition to the challenge of datasets, existing MECTEC methods fail to explicitly model the emotional and causal contexts in conversations.
% They also pay little attention to the fusion of semantic information of different levels, which leads to performance degradation. 
In addition, existing MECTEC methods fail to explicitly model emotional and causal contexts and neglect the fusion of semantic information at different levels, leading to performance degradation.
In this paper, we propose M\textsuperscript{3}HG, a novel model that explicitly captures emotional and causal contexts and effectively fuses contextual information at both inter- and intra-utterance levels via a multimodal heterogeneous graph.
% Based on the information propagated on the graph, M\textsuperscript{3}HG can accurately identify potential triplets of emotion utterance, cause utterance, and emotion category.
Extensive experiments demonstrate the effectiveness of M\textsuperscript{3}HG compared with existing state-of-the-art methods.
The codes and dataset are available at  \url{https://github.com/redifinition/M3HG}.
\end{abstract}

\section{Introduction}

Emotion Cause Analysis in Conversations (ECAC) aims at identifying emotions and their causes in conversations, which is a crucial research field in natural language processing~\cite{li2022ecpec, wang2023generative}.
% It is widely utilized for multiple analytical purposes across various social media platforms, including customer behavior analysis~\cite{jeong2019social}, mental health analysis~\cite{barros2021vader, vajre2021psychbert}, and public opinion monitoring~\cite{barreto2023sentiment, karamouzas2022public}.
However, most of ECAC research~\cite{li2022ecpec, wang2023generative, ECQED, MPEG} only focuses on the textual contexts, overlooking other modalities~\cite{soleymani2017survey}.

\begin{figure*}
\centering
  \includegraphics[width=0.9\textwidth]{Figure/introduction.pdf}
  \caption{An example of the MECTEC task. Each utterance contains three different modalities - text, audio, and video. Arrows represent causal relationships that link the cause utterances to the corresponding emotion utterances. The dashed box at the bottom lists all the <utter-cause-emotion> triplets identified in this example.}
  \label{fig:introduction figure}
  \vspace{-2mm}
\end{figure*}

To address this limitation, \citet{ecf} proposed a new task called Multimodal Emotion Cause Triplet Extraction in Conversations (MECTEC).
The task aims to simultaneously identify the emotion utterance, the corresponding cause utterances, and the emotion category (i.e., the \textit{utter-cause-emotion} triplet) from a conversation containing three modalities: text, audio, and video.
Figure~\ref{fig:introduction figure} illustrates a multimodal conversation between a mother and daughter.
In this example, there are six non-neutral utterances, and consequently, six utter-cause-emotion triplets are identified.
% Compared with the ECAC task, MECTEC is characterized by: 1) Multimodality of contexts (such as text, audio and video), which results in a more complex way of emotional expression.
% 2) Multi-scale semantic information, which comes from the overall conversation content and the topological relationships between utterances, and from details in each utterance like intonation and facial expressions.
% These characteristics present significant challenges to the development of this field.
MECTEC differs from ECAC in 1) multimodal contexts (i.e., text, audio, and video) resulting in more complex emotional expression, and 2) multi-scale semantic information from overall conversation and utterance features like intonation and facial expressions, which pose significant challenges.

Another major challenge in MECTEC is the scarcity of datasets.
% Although there have been plenty of text-based datasets for the ECAC task, only one dataset, namely the ECF dataset~\cite{ecf}, is specifically designed for MECTEC.
While numerous text-based datasets exist for ECAC, only one dataset, namely the ECF dataset~\cite{ecf}, is specifically designed for MECTEC.
% It consists of video segments with corresponding scripts and audios.
However, the videos in ECF are all from the \textit{Friends} TV series with restricted speakers and scenarios, hindering MECTEC model development.
Therefore, in this work, a new multimodal, multi-scenario MECTEC dataset, namely \textbf{MECAD}, is constructed.
% To address this, we introduce \textbf{MECAD}, a new multimodal, multi-scenario MECTEC dataset, which, to the best of our knowledge, is the first of its kind and will significantly support research in this field.
To the best of our knowledge, it is the first of its kind and will greatly facilitate research in this field.
% To the best of our knowledge, MECAD is the first MECTEC dataset with multi-scene, which will greatly facilitate the research in this field.
% To the best of our knowledge, MECAD is the first MECTEC dataset with multiple scenes, which will greatly facilitate research in this field.

Constrained by the limited dataset,  existing MECTEC models have various deficiencies.
% An early study ~\cite{ecf} proposes a two-stage architecture, which first predicts emotion utterances and cause utterances, and then pairs the emotion utterances with their cause utterances.
% However, this two-stage architecture is computationally intensive, and suffers from significant error accumulation.
% Therefore, recent studies propose one-stage architectures which employ graph neural networks or prompt engineering techniques to extract utter-cause-emotion triplets from conversations~\cite{UniMEEC, wang2023generative, li2024multimodal}.
% An early study~\cite{ecf} proposes a two-stage architecture that predicts emotion and cause utterances separately, then pairs them. 
% However, this approach is computationally intensive and prone to error accumulation.
~\citet{ecf} proposed a two-stage architecture that predicts emotion and cause utterances separately.
However, this approach is computationally intensive and prone to error accumulation.
Therefore, recent studies ~\cite{UniMEEC, wang2023generative, li2024multimodal} propose one-stage architectures using graph neural networks or prompt engineering to extract utter-cause-emotion triplets.
However, these methods do not explicitly extract specific contexts related to emotions and their causes.
% However, these methods do not explicitly extract contexts related to emotions and causes.
According to emotion attribution theory~\cite{weiner1985attributional}, the relationships of emotions and their causes are revealed by specific contexts, such as emotional words in texts, and intonations in audio and video conversations.
% For example, in Utterance 3 in Figure~\ref{fig:introduction figure}, the exclamation “feel like ages” in the text and the pleasant facial expression in the video reflect the speaker's happiness, while the event “I had your cooking” and the speaker's happy tone reflect the emotional cause of utterance 3.
For example, in Utterance 3 in Figure~\ref{fig:introduction figure},  a pleasant facial expression indicates happiness, while “haven't had your cooking” and a happy tone reveal the cause.
The example illustrates that emotions and their causes depend on contextual cues across multiple modalities, highlighting the necessity of explicitly modeling their specific contexts.

% 这里是新加的部分
% 除此之外，先前的工作均无法有效识别对话中原因话语出现在情绪话语之后的情形。比如图1中的话语1，仅通过话语1的上下文无法得到Luo开心的原因，这需要分析对话的整体内容才能得知。
In addition, previous work~\cite{ecf, UniMEEC, wang2023generative, li2024multimodal, wei2020effective} fail to effectively identify the \textbf{cause utterances occurring after emotion utterances}. 
% For example, in Utterance 1 in Figure 1, the reason why Luo is happy cannot be obtained from the context of Utterance 1 alone, which needs to be analyzed in the context of the conversation as a whole.
For example, in Utterance 1 in Figure 1, the reason why Luo is happy cannot be obtained only from the historical context of Utterance 1.
To find out the real cause of emotion in Utterance 1, the whole conversation should be scrutinized, which is overlooked by previous work.
% 现有模型缺乏处理这类情形的有效改进
% Existing MECTEC models lack effective solutions for this scenario.

% Furthermore, existing models~\cite{ecf, UniMEEC, wang2023generative, li2024multimodal} perform multimodal fusion in separate stages, leading to the loss of contextual information.
Furthermore, existing models~\cite{ecf, UniMEEC, wang2023generative, li2024multimodal} fail to adequately extract semantic information at different scales.
% As shown in Figure~\ref{fig:introduction figure}, the semantic information which reveals the relationship of an utterance and its cause resides not only in connections between different utterances but also in the content of each utterance.
As shown in Figure~\ref{fig:introduction figure}, the semantic information that reveals the relationship of an utterance and its cause not only resides in inter-connections between utterances but also resides in the intra-content of each utterance.
% As shown in Figure~\ref{fig:introduction figure}, the semantic information which reveals the relationship of an utterance and its cause resides in connections not only between different utterances but also different modalities of each utterance.
Therefore, it's essential to comprehensively integrate semantic information in different scales during modality fusion.

To solve the aforementioned problems, we propose an MECTEC model based on the multimodal, multi-scale, and multi-type node heterogeneous graph, named \textbf{M\textsuperscript{3}HG}.
M\textsuperscript{3}HG accurately extracts emotion and cause-related contexts and fuses multimodal, multi-scale semantic information using multimodal heterogeneous graph attention network (HGAT) with multi-type nodes.

Our contributions can be summarized as follows:
\begin{itemize}[itemsep=0pt,parsep=0pt,topsep=0pt,partopsep=0pt]
  \item The first Chinese multi-scenario MECTEC dataset, \textbf{MECAD}, and an online sentiment data annotation toolkit are constructed.
  The dataset consists of 989 conversations with 10,519 utterances annotated with important information such as emotion labels, their causes, and types of emotional causes.
  It will greatly benefit the development of models in the MECTEC and related fields.
  \item An efficient MECTEC model, namely M\textsuperscript{3}HG, is proposed to identify utter-cause-emotion triplets from multimodal conversations.
  % M\textsuperscript{3}HG explicitly extract emotion- and cause-related contextual features, thus enhancing the association between emotions and causes.
  It explicitly extracts specific emotion and cause-related contexts to find connections between emotions and causes.
  % Besides, it fully integrates semantic information from inter and intro-utterance scales and effectively enhances the model's ability to handle situations in which cause utterances appear after emotion utterances.
  Besides, it fully integrates semantic information from inter and intra-utterance levels to enhance the model's predictive ability. 
  \item Extensive experiments are performed to verify the performance of our proposed model and other state-of-the-art models on MECAD and ECF datasets.
  Experimental results reveal that M\textsuperscript{3}HG outperforms its counterparts, which demonstrates the effectiveness of our model.
\end{itemize}
% \begin{enumerate}
%     \item[(1)] We develop the first Chinese multi-scenario MECTEC dataset, \textbf{MECAD}, along with an online sentiment data annotation toolkit to support the development of models in the MECTEC and related fields.
%     Additionally, a comprehensive data cleaning and assessing process is performed on MECAD to ensure its data quality.
%     \item[(2)] We propose an efficient MECTEC model, namely \textbf{M\textsuperscript{3}HG}, to extract utter-cause-emotion triplets from conversations.
%     M\textsuperscript{3}HG utilizes a novel multimodal graph attention network (GAT) to explicitly extract emotion- and cause-related contextual features, thus enhancing the association between emotions and causes.
%     In addition, it fully integrates semantic information from different scales, thus extracting the utter-cause-emotion triplet more comprehensively and accurately.
%     \item[(3)] We perform extensive experiments to verify the performance of our proposed model as well as other state-of-the-art models on MECAD and ECF datasets.
%     The experimental results reveal that M\textsuperscript{3}HG outperforms the state-of-the-art methods on both datasets, which demonstrates the effectiveness of our model.
% \end{enumerate}

\section{Related Works}
\textbf{Emotion Cause Analysis in Conversations.}
Most existing studies on ECAC focus on Causal Emotion Entailment (CEE) and Emotion Cause Pair Extraction in Conversations (ECPEC).
CEE aims to identify which cause utterances trigger the non-neutral emotions of the target utterances.
Since CEE assumes emotion utterances are given, most related work~\cite{RECCON, li2022neutral, zhang2022tsam, gu2023page} viewed CEE as an utterance classification problem.
However, because emotions of utterances are often unknown in real-world conversations, \citet{li2022ecpec} proposed the ECPEC task which additionally predicts emotions for the target utterances.
Subsequent work~\cite{wang2023generative, zhao2023knowledge} has incorporated commonsense knowledge into GATs to improve the model's semantic understanding of emotions and causes, achieving better performance.
Besides, some methods~\cite{ECPE-2D, ECPE-MLL, wei2020effective, UECA-Prompt} from models in the Emotion cause Pair Extraction (ECPE) field are also adapted for the ECPEC task.

\noindent \textbf{Multimodal Emotion Cause Triplet Extraction in Conversations.} In recent years, multimodal conversation scenarios on social media platforms have grown significantly, as more individuals share their lives and express emotions through live streaming and various online chats.
To advance emotion cause analysis in multimodal conversation scenarios, \citet{ecf} introduced the MECTEC task and released the ECF dataset.
% Since this task was recently introduced, there are few solutions provided.
However, few solutions have been proposed for this recently introduced task.
\citet{li2024multimodal} incorporated emotion transition information into emotion-cause pair extraction using a novel labeling constraint, while \citet{UniMEEC} fused semantic information across modalities via prompt engineering.
These methods treat multimodal fusion and contextual information extraction for emotional causes as separate processes.
Furthermore, they fail to effectively integrate semantic information across different scales, which significantly hampers the overall performance of models in the MECTEC task.
% To overcome these limitations, we propose an MECTEC model that fully integrates multi-scale semantic information across different modalities. This approach effectively prevents the loss of contextual information during modality fusion, thereby enhancing the triplet extraction accuracy for the MECTEC task.
To address these issues, we propose a model that fully integrates multi-scale semantic information from different modalities, preventing the loss of contextual information during fusion and improving triplet extraction accuracy.

\noindent\textbf{Datasets for the ECAC Task.}
\begin{table}
  \small
  \caption{A summary of datasets for ECAC task. T, A, V stand for text, audio and video respectively.}
  \label{tab:related-dataset}
  \centering
  \begin{tabular}{l|c|c|c}
    \toprule
    Dataset & Modalities & Sources &  \# Instances\\
    \midrule
    % ECE Corpus~\cite{ECE-Corpus} & T & News & 2,105\\
    % Weibo-Emotion~\cite{Weibo-Emotion} & T & Microbog  &  7,000 \\ 
    % NTCIR-13-ECA~\cite{NTCIR-13-ECA} & T & Fiction  & 2,403 \\
    % REMAN~\cite{REMAN} & T & News  & 1,720 \\
    % GoodNewsEveryone~\cite{GoodnewsEveryone} & T & News  & 5,000 \\
    % \midrule
    % RECCON~\cite{RECCON} & T & Act and Daily  & 11,769 \\
    RECCON & T & Act and Daily  & 11,769 \\
    % ConvECPE~\cite{li2022ecpec} & T & Act  & 7,433 \\
    ConvECPE & T & Act  & 7,433 \\
    \midrule
    % ECF~\cite{ecf} & T,A,V & TV \textit{Friends}  & 13,509 \\
    ECF & T,A,V & TV \textit{Friends}  & 13,509 \\
    % \textbf{MECAD} (ours) & T,A,V & 56 TV series &  10,516 \\ 
    \textbf{MECAD} & T,A,V & 56 TV series &  10,516 \\
  \bottomrule
  \end{tabular}
  \vspace{-4mm}
\end{table}
Table \ref{tab:related-dataset} summarizes popular datasets in ECAC.
% ~\citet{RECCON} firstly introduced the RECCON dataset for the ECAC task.
% Additionally, \citet{li2022ecpec} built the ConvECPE dataset based on the RECCON dataset.
\citet{RECCON} introduced the RECCON dataset for the ECAC task, and \citet{li2022ecpec} extended it by building the ConvECPE dataset.
% Considering the multimodal nature of conversations, \citet{ecf} recently developed the ECF dataset for MECTEC.
Given the multimodal nature of conversations, \citet{ecf} developed the ECF dataset for MECTEC.
% However, all the scenes in this dataset are from the \textit{Friends} TV series, bringing in the issue of limited conversation scenarios and similar conversation contents.
However, all scenes in ECF are drawn from the \textit{Friends}, limiting the diversity of conversation scenarios and contents.

\section{Proposed MECAD Dataset}
To facilitate the research in MECTEC and other related fields, we constructed a multi-scenario MECTEC dataset called MECAD.
Compared with ECF~\cite{ecf}, MECAD has more diverse conversation scenarios.
% In addition to the emotion categoty of each utterance and its cause, we also labeled the types of emotion causes (such as \textit{event}, \textit{expression}, etc) and the modal basis of annotation (i.e., the cause is annotated based on \textit{text}, \textit{audio}, or \textit{video}) for future multimodal emotion cause analysis studies.
In addition to labeling emotion categories and their causes for each utterance, we also categorized the types of emotion causes (e.g., \textit{event}, \textit{expression}) and the modality of annotation (i.e., \textit{text}, \textit{audio}, or \textit{video}) to support future studies in multimodal emotion cause analysis.
% We labeled emotion categories,the corresponding cause utterances, types of emotion causes (e.g., \textit{event}, \textit{expression}, etc), and the modal basis of annotation (\textit{text}, \textit{audio}, or \textit{video}).

We selected the publicly available M\textsuperscript{3}ED~\cite{M3ED} dataset as our data source, which contains 990 segments from 56 Chinese TV series.
% These segments provide the spontaneous emotional interactions between users in social media~\cite{wirth2005media}.
% It is believed that media contents contain the spontaneous emotional interactions between users~\cite{wirth2005media}.
However, M\textsuperscript{3}ED dataset only contains conversation scripts, audios, and screenshots, lacking corresponding videos.
Therefore, we endeavored to collect the corresponding video segments based on the conversation timestamps provided by M\textsuperscript{3}ED.
% We concatenated incomplete sentences to form 10,516 full utterances across 989 multimodal conversations.
We concatenated sentences to form 989 multimodal conversations with 10,516 full utterances.

We invited 10 Chinese graduate students majored in Psychology to annotate the corresponding cause utterances, the types of emotion causes and the modal cues of annotations in the conversations. 
To obtain high-quality annotations, we designed detailed guidelines based on previous studies~\cite{dirven1997emotions, steptoe2009emotional}, trained the volunteers, and tested them with annotation cases.
Only those passing the test participated in the final annotation process.
Each volunteer was paid \$50 for their annotations.
% Then, we randomly assigned three qualified annotators for each conversation, and the annotators were required to label in sequence the corresponding cause utterance, the type of emotion cause, and the modal basis of the labeling for each non-neutral utterance.
Then, we randomly assigned three qualified annotators for each conversation.
If divergence exists among annotations from different volunteers, the final annotation for the utterance is determined by majority voting.
Two strategies were used to review and revise incorrect annotations:
% 1) We calculated annotation consistency among the three annotators for each TV series.
1) Annotation consistency among the three annotators for each TV series is calculated.
For series with low consistency, the annotators rechecked and revised their labels as needed.
% 2) If disagreements remained, a fourth annotator was invited to relabel the conflicting utterances and make the final decision.
2) If disagreements remained, a fourth annotator was invited to relabel the utterances and make the final decision.

To enhance annotation efficiency and accuracy, we developed an online multimodal conversation emotion cause annotation tool.
% This is a generalized toolkit for multimodal annotation of conversations, with features including quick configuration of annotation content, multi-user and task management, annotation consistency management, and distributable deployment. 
The interface of the annotation tool is shown in Figure~\ref{fig:annotation toolkit} in Appendix~\ref{Appendix-toolkit}.
% This annotation tool can be used for work related to sentiment analysis.
This tool is highly reusable and user-friendly, making it ideal for related research in the future.

% To comprehensively evaluate the quality and consistency of annotations, we employ Cohen’s Kappa~\cite{cohens-kappa} to assess the agreement between two annotators, and Fleiss’s Kappa~\cite{fless-kappa} to measure the overall consistency among all annotators.
% The results of the Cohen’s Kappa score are shown in Appendix~\ref{appendix-b}.
% In addition, the Fleiss’s Kappa score of 0.6932 exceeds the threshold of 0.61~\cite{landis1977measurement}, which indicates that our labeling results are statistically significant.
We use Cohen’s Kappa~\cite{cohens-kappa} to assess pairwise agreement and Fleiss’s Kappa~\cite{fless-kappa} for overall consistency among annotators.
The Cohen’s Kappa results are in Appendix~\ref{appendix-b}, and the Fleiss’s Kappa score of 0.6932 exceeds the threshold of 0.61~\cite{landis1977measurement}, confirming the statistical reliability of our annotations.

% The dataset statistics and detailed analysis of MECAD are presented in Figure~\ref{fig:annotation result} in Appendix~\ref{appendix-a}.
% MECAD provides solid support for assessing the performance and generalization capabilities of MECTEC models in broader conversation scenarios.
The dataset statistics and detailed analysis of MECAD are presented in Figure~\ref{fig:annotation result} in Appendix~\ref{appendix-a}.
% demonstrating its utility for evaluating MECTEC models’ performance and generalization in broader conversational scenarios.
MECAD provides solid support for assessing the performance and generalization capabilities of MECTEC models in broader scenarios.
\section{\texorpdfstring{Framework of Proposed M\textsuperscript{3}HG}{Framework of Proposed M3HG}}

\begin{figure*}
\centering
  \includegraphics[width=1\textwidth]{Figure/OverallFramework.pdf}
  %\caption{The overview of our proposed multimodal, multi-scale, and multi-type node heterogeneous graph for MECTEC (M\textsuperscript{3}HG).}
  \caption{The framework of proposed M\textsuperscript{3}HG. It consists of three main components: unimodal feature extraction, graph construction and multi-scale semantic fusion, and emotion-cause classification.}
  \label{fig:overallframework}
\end{figure*}

\subsection{Task Definition}
% Given a conversation $C = \{(S_1, U_1),\!\dots,\!(S_i, U_i),\!\dots,\!(S_n, U_n)\}$, where $S_i$ denotes the speaker of the $i$-th utterance $U_i$, $n$ denotes the length of the conversation $C$, $U_i = \{U_i^t, U_i^a, U_i^v\}, i \in \{1, 2, \dots, n\}$, and $t$, $a$, $v$ are the text, audio and video modality, respectively.
Given a conversation $C = \{(S_i, \bm{U}_i)\}_{1\leq i\leq n}$, where $S_i$ denotes the speaker of the $i$-th utterance $\bm{U}_i$, $n$ denotes the length of the conversation $C$, $\bm{U}_i = \{U_i^t, U_i^a, U_i^v\}$, and $t$, $a$, $v$ are the text, audio and video modality, respectively.
The goal of MECTEC is to identify all the utter-cause-emotion triplets from the conversation $C$:
\begin{equation}
% \mathcal{P} = \{\dots, (U_j^e, U_j^c, y_j^e), \dots\},
\mathcal{P} = \{(\bm{U}_j^e, \bm{U}_j^c, y_j^e)\},
\end{equation}
where $\bm{U}_j^e$ is the $j$-th utterance with emotion $y_j^e$, $\bm{U}_j^c$ is the corresponding cause utterance, and $y_j^e\in \{\textit{Anger}, \textit{Disgust}, \textit{Fear}, \textit{Joy}, \textit{Sadness}, \textit{Surprise}\}$~\cite{ekman1992argument}.
% The emotion category $y_j^e$ is one of six non-neutral emotions~\cite{ekman1992argument}, i.e., $y_j^e\in \{\textit{Anger}, \textit{Disgust}, \textit{Fear}, \textit{Joy}, \textit{Sadness}, \textit{Surprise}\}$.

\subsection{Model Overview}
M\textsuperscript{3}HG is an end-to-end (E2E) MECTEC model, as illustrated in Figure~\ref{fig:overallframework}.
It consists of four key components: \textit{unimodal feature extraction}, \textit{graph construction}, \textit{multi-scale semantic fusion}, and \textit{emotion-cause classification}.

In \textit{unimodal feature extraction},  M\textsuperscript{3}HG extracts local contextual representations for each utterance using modality-specific feature extractors and unimodal encoders.
In \textit{graph construction}, M\textsuperscript{3}HG constructs a conversation interaction graph using these feature representations to explicitly model the emotion and cause-related contexts.
In \textit{multi-scale semantic fusion}, M\textsuperscript{3}HG combines semantic information at different scales within the conversation interaction graph to produce a comprehensive feature representation of both emotion and cause contexts.
% Finally, in emotion-cause classification, emotion-cause pairs are predicted by a link predictor that fully utilizes emotion and cause contextual representations.
In \textit{emotion-cause classification}, emotion and cause contextual representations are concatenated and used to extract utter-cause-emotion triplets with with position embedding.

\subsection{Unimodal Feature Extraction}
% First, we utilize SA-RoBERTa~\cite{SA-RoBERTa}, Wav2Vec2~\cite{wav2vec2}, and DenseNet~\cite{DenseNet} to extract three feature representations $\bm{E}^t$, $\bm{E}^a$, and  $\bm{E}^v$, from text, audio, and video, respectively, where $\bm{E}^t \in \mathbb{R}^{n \times d_t}$, $\bm{E}^a \in \mathbb{R}^{n \times d_a}$, and  $\bm{E}^v \in \mathbb{R}^{n \times d_v}$, and  $d_t$, $d_a$, $d_v$ represent dimensions of the hidden layer representations of the three modality-specific feature extraction models.
First, we utilize SA-RoBERTa~\cite{SA-RoBERTa}, Wav2Vec2~\cite{wav2vec2}, and DenseNet~\cite{DenseNet} to extract three feature representations $\bm{E}^t$, $\bm{E}^a$, and  $\bm{E}^v$, from text, audio, and video, respectively, where $\bm{E}^t \in \mathbb{R}^{n \times d_t}$, $\bm{E}^a \in \mathbb{R}^{n \times d_a}$, and  $\bm{E}^v \in \mathbb{R}^{n \times d_v}$, and  $d_t$, $d_a$, $d_v$ represent dimensions of the hidden layer representations of the three modalities.
The extraction process is described in Appendix~\ref{appendix-c.1}.

Then, we encode each feature representation within an unimodal local context.
% For text, we further encode $\bm{E}^t$ using the masked multi-head self-attention mechanism (MHSA)~\cite{self-attention}, which captures the local contextual information through the mask $\bm{M}$ and the utterance window $w$. 
% The mask is set to 0 within the utterance window and to negative infinity outside the window.
% This can be expressed as follows:
% \begin{equation}
% \bm{M}[i,j] = 
% \begin{cases} 
% 0, & \text{if } |T_i - T_j| \leq w, \\
% -\infty, & \text{otherwise},
% \end{cases}
% \end{equation}
% where $i$ and $j$ are the marker positions in the input sequence $i$, and $T_i$ and $T_j$ are the indices of the utterances where the $i$-th and $j$-th markers are located.
% In the MHSA block, we apply the mask $\bm{M}$ to $\bm{E}^t$, and it can be expressed as follows:
% \begin{equation}
% \bm{H}^t = MHSA(\bm{E}^t, \bm{M}).
% \end{equation}
For text, we apply multi-head self-attention~\cite{self-attention} to $\bm{E}^t$ to capture local contextual information, resulting in $\bm{H}^t$.
For $\bm{E}^a$ and $\bm{E}^v$, we use a GRU-based network~\cite{CFN-ESA} to extract local context by leveraging the RNN structure’s capability to handle temporal features, which is expressed as:
\begin{equation}
\begin{aligned}
\bm{E}'^m &= LN(\bm{E}^m + GRU(\bm{E}^m)), \\
\bm{H}^m &= LN(\bm{E}^m + \bm{E}'^m + FFN(\bm{E}'^m),
\end{aligned}
\end{equation}
where $\bm{H}^m \in \mathbb{R}^{n \times d_m},m \in \{a,v\}$, $LN$ denotes layer normalization, and $FFN$ denotes a feedforward neural network.

After encoding the local context for each modality, we obtain the sequence representations $\bm{H}^t$, $\bm{H}^a$, $\bm{H}^v$ for text, audio, and video.
We then apply three linear layers to map $\bm{H}^t$, $\bm{H}^a$, $\bm{H}^v$ to $\bm{H}'^t$, $\bm{H}'^a$, $\bm{H}'^v$ with the same dimension $d_h$.

\subsection{Graph Construction}
% In order to enable M\textsuperscript{3}HG to effectively fuse semantic information of different scales from different modalities, we need to construct a heterogeneous graph that can effectively represent the connections between/within utterances and the interactions between different modalities.
To enable M\textsuperscript{3}HG to fuse multi-scale semantic information across modalities, we construct a heterogeneous graph that represents both inter- and intra-utterance connections, as well as cross-modal interactions.
The structure of this heterogeneous graph can be denoted by $\mathcal{G}=(\mathcal{V},\mathcal{E},\mathcal{R})$, where $\mathcal{V}$ is the node set consisting of all graph nodes $v_i$, $\mathcal{R}$ is the relation set consisting of all relations $r_{ij}$ between any two nodes $v_i$ and $v_j$, and $\mathcal{E}$ is the edge set consisting of all edges represented as $(v_i,r_{ij},v_j)$.
% The graph is constructed as follows.


\noindent\textbf{Nodes.} To explicitly model emotion and cause-related contexts in conversations, we model them as \textit{emotional context nodes} $N^e$ and \textit{causal context nodes} $N^c$, respectively.
To enable $\mathcal{G}$ to accurately perceive the conversation information, we model the whole conversation as a \textit{conversation node}.
% In addition, we model each utterance using the \textit{utterance node}.
% Since each utterance contains three modal features, both the utterance node and the conversation node are specially designed as Super-Nodes which contain three modal features and are denoted as $SN^u$ and $SN^d$, respectively.
Each utterance is represented by an \textit{utterance node}. 
% Since each contains three modal features, both the utterance node and conversation node are designed as Super-Nodes containing these three modalities, denoted as $SN^u$ and $SN^d$.
Both the utterance node and conversation node are designed as Super-Nodes containing these modalities, denoted as $SN^u$ and $SN^d$, since they contain three modal features.
Therefore, $\mathcal{G}$ contains four types of nodes: $N^e$, $N^c$, $SN^u$ and $SN^d$.

$N^e$ and $N^c$ are first initialized with textual sequence representations $\bm{H}'^t$, then updated with contextual information from the other two modalities, which is described in Section~\ref{method-4.5}.
% Each utterance Super-Node $SN^u$ consists of three nodes, i.e., $SN^u=\{N^t,N^a,N^v\}$, which are initialized using the local context sequence representation of the three modalities, i.e., $\bm{H}'^t,\bm{H}'^a,\bm{H}'^v$.
Each utterance Super-Node $SN^u =\{N^t,N^a,N^v\}$ is initialized using $\bm{H}'^t,\bm{H}'^a,\bm{H}'^v$.
% The conversation node $SN^d$ also consists of three nodes corresponding to three modalities, i.e., $SN^d=\{N_d^t,N_d^a,N_d^v\}$. 
% Since the conversation node needs to fully contain the global conversation information, it is initialized using the average pooling results of $\bm{H}'^t,\bm{H}'^a,\bm{H}'^v$.
% The conversation node $SN^d=\{N_d^t,N_d^a,N_d^v\}$ is initialized by averaging $\bm{H}'^t,\bm{H}'^a,\bm{H}'^v$ to capture global conversation information.
The conversation node $SN^d=\{N_d^t,N_d^a,N_d^v\}$ is initialized by averaging $\bm{H}'^t,\bm{H}'^a,\bm{H}'^v$ to capture global information.

\noindent\textbf{Edges and Relations.} There are five types of Super-Edges connecting the aforementioned Super-Nodes: \textit{same speaker} ($r_{ss}$), \textit{different speaker} ($r_{ds}$), \textit{global connection} ($r_{gc}$), \textit{emotion connection} ($r_{ec}$) and \textit{cause connection} ($r_{cc}$).
The \textit{same speaker} edge connects the utterance Super-Nodes $SN^u$ from the same speaker.
Inspired by the work of \citet{DAG-ERC}, we define the local context as $K$ preceding utterances from the same speaker of $SN^u$, where $K$ is a hyper-parameter.
% The utterance Super-Nodes in local context from the same speaker are connected to $SN^u$ by directed \textit{same speaker} edges.
The \textit{different speaker} edge connects the utterance Super-Nodes within the local context from different speakers to $SN^u$.
The bidirectional \textit{global connection} edge connects all the utterance Super-Nodes $SN^u$s with the conversation Super-Node $SN^d$, facilitating the propagation of global contextual information.
The \textit{emotion connection} edge and the \textit{cause connection} edge 
connect $SN^u$ with its corresponding emotional context node $N^e$ and causal context node $N^c$, respectively.
They explicitly capture the emotion and cause context specific to each utterance.

% It is worth mentioning that M\textsuperscript{3}HG is the first MECTEC model that can explicitly and efficiently deal with the situation that \textbf{cause utterances appear after emotion utterances} in a conversation, since each utterance can be messaged through the \textit{global connection} node.
M\textsuperscript{3}HG is the first MECTEC model capable of handling situations where \textbf{cause utterances appear after emotion utterances}, as each utterance is linked through the \textit{global connection} node.
% 附录中详细的实验进一步说明了这一点。
The detailed experiments in Appendix~\ref{in-depth analysis} further validate this capability.
The pseudo-code of graph construction and a constructed graph for the conversation in Figure~\ref{fig:introduction figure} are provided in Appendix~\ref{appendix:pseudo-code} and Appendix~\ref{appendix-c.3}, respectively.
The graph construction process of M\textsuperscript{3}HG can be expressed as:
 % 开始缩小字体
 % \begin{small}
\begin{equation}
    \begin{aligned}
    \mathcal{G} &= (\mathcal{V}, \mathcal{E}, \mathcal{R}), \\
    \mathcal{V} &= \{ SN_i^u, N_i^e, N_i^c, SN^d \}_{1\leq i\leq n}, \\
    SN_i^u &= \{ N_i^t, N_i^a, N_i^v \},\\
    SN^d &= \{ N_d^t, N_d^a, N_d^v\}, \ \\
    \mathcal{R} &= \{ r_{ss}, r_{ds}, r_{gc}, r_{ec}, r_{cc} \}, \\
    % \mathcal{E} &= \{ \ldots, (v_i, r_{ij}, v_j), \ldots \}, \quad v_i, v_j \in \mathcal{V}, \quad r \in \mathcal{R}.
    \mathcal{E} &= \{(v_i, r_{ij}, v_j)\}, v_i, v_j \in \mathcal{V}, r \in \mathcal{R},
    \end{aligned}
\end{equation}
% \end{small}
where superscripts $u,e,c,d$ denote node types,  and $m$ denotes three modalities.
Based on the constructed graph $\mathcal{G}$, the emotion and cause contexts are effectively modeled.

\subsection{Multi-scale Semantic Information Fusion}
\label{method-4.5}
Based on graph $\mathcal{G}$, we designed a comprehensive approach to integrate semantic information across different modalities and scales.
% Our multi-scale semantic information fusion mechanism contains two levels, i.e., inter-utterance-level fusion and intra-utterance-level fusion.
% The inter-utterance-level fusion propagates semantic information across different modalities among utterances as well as the conversation contexts.
% The intra-utterance-level fusion propagates semantic information and emotion cause-related contexts inside the utterance.
% This mechanism operates at two levels: inter-utterance fusion, which propagates semantic information among utterances and conversation-level contexts, and intra-utterance fusion, which captures emotion- and cause-related cues within individual utterances.
This mechanism is implemented in two levels: \textbf{intra-utterance fusion} which captures emotion and cause-related contexts within utterances, and \textbf{inter-utterance fusion} which propagates semantic information among utterances and conversation-level contexts.
Both levels leverage HGAT~\cite{HAN} to propagate and fuse semantic information through various meta-paths~\cite{HAN} within $\mathcal{G}$.
% This allows node features to be thoroughly updated by incorporating multi-scale semantic information.
This ensures thorough updates to node features by integrating multi-scale semantic information.
% This enables the thorough update of node features by taking semantic information from multiple scales into account.

The meta-paths in $\mathcal{G}$ are defined as:
% \begin{small}
\begin{equation}
\begin{aligned}
\Phi = \{\phi(v_i, r_{ij}, v_j)\}, v_i, v_j \in \mathcal{V}, \\
\phi(v_i, r_{ij}, v_j) = v_i \xleftrightarrow{r_{ij}} v_j, r_{ij} \in \mathcal{R},
\end{aligned}
\end{equation}
% \end{small}
where $\phi(v_i, r_{ij}, v_j)$ represents all paths that connect node $v_i$ to node $v_j$ via edge type $r_{ij}$.

\noindent\textbf{Intra-utterance-level Fusion.} As shown in Figure~\ref{fig:overallframework}, for each utterance Super-Node $SN^u$, we perform intra-utterance-level fusion by integrating semantic information within the utterance.
We define the meta-path $\Phi_{intra}$ for  intra-utterance-level semantic fusion for $SN_n^u$ as:
\begin{equation}
\begin{aligned}
% \Phi_{intra} = &\left\{\ldots, \phi\left(N^{m_1}, N^{m_2}, r_{m_1,m_2}\right), \ldots\right\} \\
\Phi_{intra} = &\{\phi\left(N^{m_1}, N^{m_2}, r_{m_1,m_2}\right)\}\\
% &\cup \left\{\ldots, \phi\left(N^m, N^e, r_{m,e}\right), \ldots\right\} \\
&\cup \{\phi\left(N^m, N^e, r_{m,e}\right)\} \\
% &\cup \left\{\ldots, \phi\left(N^m, N^c, r_{m,c}\right), \ldots\right\},
&\cup \{\phi\left(N^m, N^c, r_{m,c}\right)\},
\end{aligned}
\end{equation}
where $m_1, m_2, m \in \{t, a, v\}$, $N^m$ represents the nodes of modality $m$ within the $SN^u$, and $r_{m_1,m_2}$ denotes the edges connecting $N^{m_1}$ and $N^{m_2}$. 
$r_{m,e}$ denotes edges connecting nodes $N^m$ to the emotional context nodes $N^e$, facilitating the aggregation of emotional contexts conveyed by different modalities within the utterance.
Similarly, $r_{m,c}$ represents the edges that connect $N^m$ to the causal context nodes $N^c$, enabling the aggregation of causal contexts.
$\Phi_{intra}$ effectively models the process of semantic information fusion in a single utterance.

Next, we incorporate node-level attention into $\Phi_{intra}$.
For each meta-path in $\Phi_{intra}$ and nodes $v_i\in \{N^m,N^e,N^c\}$, the importance of its neighbors $\mathcal{N}_i$ in $\Phi_{intra}$ is computed as:

\begin{small}
\begin{equation}
\alpha_{ij}^{\phi} = \frac{\exp\left(\sigma\left(\bm{a}_\phi^T \cdot \left[\bm{H}'_i \parallel \bm{H}'_j\right]\right)\right)}{\sum_{k \in \mathcal{N}_i^{\phi}} \exp\left(\sigma\left(\bm{a}_\phi^T \cdot \left[\bm{H}'_i \parallel \bm{H}'_k\right]\right)\right)}, \quad \phi \in \Phi_{intra},
\label{eq:11}
\end{equation}
\end{small}
where $\sigma$ denotes the activation function, and $\bm{a}_\phi$ is the node-level attention vector of meta-path $\phi$.
% Then, the representation of node $v_i$ based on meta-path $\phi$ is obtained by aggregation through the following process:
The node representation of $v_i$ based on meta-path $\phi$ is obtained by:
\begin{equation}
% \bm{Z}_i^{\phi} = \sigma( \sum_{j \in \mathcal{N}_i^{\phi}} \alpha_{ij}^{\phi} \cdot \bm{H}'_j ).
\bm{Z}_i = \sigma( \sum_{j \in \mathcal{N}_i^{\phi}} \alpha_{ij}^{\phi} \cdot \bm{H}'_j ).
\label{eq:12}
\end{equation}

% Based on the above process, we can get the contextual features $\bm{Z}_i^\phi\in\mathbb{R}^{{1\times d}_h}$ of nodes $v_i$ under the intra-utterance-level meta-paths $\Phi_{intra}$.
This process yields the contextual features $\bm{Z}_i \in \mathbb{R}^{1 \times d_h}$ for nodes $v_i$ under the intra-utterance-level meta-paths $\Phi_{intra}$.

\noindent\textbf{Inter-utterance-level Fusion.} As illustrated in Figure~\ref{fig:overallframework}, for any two utterance Super-Nodes $SN_i^u$ and $SN_j^u$ in $\mathcal{G}$, along with the conversation Super-Node $SN^d$, we perform inter-utterance-level fusion by connecting $SN_i^u$ and $SN_j^u$ to $SN^d$, thereby integrating contextual information across utterances.
We define meta-paths $\Phi_{inter}$ for inter-utterance-level fusion between $SN^u$ and $SN^d$:
% enabling effective semantic integration across utterances:
\begin{equation}
\begin{aligned}
% \Phi_{inter} = &\{\ldots, \phi (N_i^{m_1}, N_j^{m_2}, r_{m_1,m_2}), \ldots\} \\
% &\cup \{\ldots, \phi(N_i^m, N_d^m, r_{d,m}), \ldots\} \\
% &\cup \{\ldots, \phi(N_j^m, N_d^m, r_{d,m}), \ldots\},
\Phi_{inter} = &\{\phi (N_i^{m_1}, N_j^{m_2}, r_{m_1,m_2})\} \\
&\cup \{ \phi(N_i^m, N_d^m, r_{d,m})\} \\
&\cup \{ \phi(N_j^m, N_d^m, r_{d,m})\},
\end{aligned}
\end{equation}
where $m_1, m_2, m \in \{t, a, v\}$, $N_i^m$ and $N_j^m$ represent the nodes of modality $m$ inside $SN_i^u$ and $SN_j^u$, respectively, $r_{m_1,m_2}$ denotes the edges connecting $N_i^{m_1}$ and $N_j^{m_2}$, and $r_{d,m}$ represents the edges connecting $SN^u$s to $SN^d$ in modality $m$.

The utterance information from each modality can be passed to $SN^d$ though $\Phi_{inter}$, which accomplishes inter-utterance-level fusion between utterances.
As a result, $SN^d$ comprehensively integrates information across all three modalities.
The meta-path set $\Phi_{inter}$ models multimodal connections between utterances, enabling conversation information aggregated in $\mathcal{G}$.

% Using Eqs. \ref{eq:11}-\ref{eq:12}, contextual representations of $SN^u$ and $SN^d$ are derived under the meta-path $\Phi_{inter}$ using the node-level attention mechanism.
Similar to Eq.~\ref{eq:11} and Eq.~\ref{eq:12}, the contextual representations of $SN^u$ and $SN^d$ are obtained under the meta-path $\Phi_{inter}$ by the node-level attention block.

% After completing the multi-scale semantic fusion using $\Phi_{intra}$ and $\Phi_{inter}$, we apply the semantic attention mechanism \cite{HAN} to perform meta-path-level feature fusion to incorporate multi-scale semantic information from three modalities.
% Inspired by \cite{MPEG}, after each round of multi-scale semantic information fusion, we apply a non-linear transformation to the node features using a position-wise feed-forward network (PFFN) layer.
% The process can be expressed as follows:
% \begin{equation}
% \bm{H}_{i}^{\prime \tau+1} = LN(Dropout(\bm{W}_2 \sigma(\bm{W}_1 \bm{H}_i^{\prime \tau})) + \bm{H}_i^{\prime \tau}),
% \end{equation}
% where $\bm{H}_i^{\prime \tau}$ and $\bm{H}_i^{\prime \tau+1}$ are the node embeddings of node $v_i$ at $\tau$-th iteration and ($\tau$+1)-th iteration, i.e., $\bm{H}_i^{\prime \tau} = \bm{Z}_i^{ \tau}$, $\bm{H}_i^{\prime \tau+1}=\bm{Z}_i^{\tau+1}$.
% $\bm{W}_1$, $\bm{W}_2\in\mathbb{R}^{d_h\times1}$ denote the weight matrix of the two-layer 1D-convolution, and $\sigma$ denotes the ELU activation function.
% The emotional context node  representation $\bm{Z}^e$ and the causal context node feature representation $\bm{Z}^c$ can be obtained at the end of iterations of the multi-scale semantic fusion and PFFN layers. 
After performing multi-scale semantic fusion with $\Phi_{intra}$ and $\Phi_{inter}$, we apply the semantic attention mechanism~\cite{HAN} to each node embedding $\bm{Z}_i$, integrating multi-scale semantic information from all three modalities.
Following \cite{MPEG}, each fusion iteration is followed by a position-wise feed-forward network (PFFN) layer, which updates node features through a non-linear transformation.
The emotional context node  representation $\bm{Z}_i^e$ and the causal context node feature representation $\bm{Z}_i^c$ can be obtained at the end of iterations of the multi-scale semantic fusion and PFFN layers.

\subsection{Emotion-cause classification}
For each utterance $\bm{U}_i$, its $\bm{Z}_i^e$ and $\bm{Z}_i^c$ are fed into the emotion-specific Multi-Layer Perceptron (Emotion MLP) and the cause-specific Multi-Layer Perceptron (Cause MLP) to predict its emotion category $\hat{y}_i^e$ and the cause indicator $\hat{y}_i^c$ which indicates whether $\bm{U}_i$ can be a cause utterance.
For each utterance pair $\bm{U}_i$ and $\bm{U}_j$, we compute a relative position encoding $RPE_{ij}$ to capture the positional relationship between $\bm{U}_i$ and $\bm{U}_j$.
% Following RankCP~\cite{wei2020effective}, we utilize the RBF kernel function to compute $RPE_{ij}$.
We utilize the RBF kernel function~\cite{wei2020effective} to compute $RPE_{ij}$, which captures the relative positional relationships between utterances through a nonlinear relation.
$\bm{Z}_j^e$, $\bm{Z}_i^c$ and $RPE_{ij}$ are then concatenated and fed into a new MLP to determine whether $\bm{U}_i$ is the cause utterance of $\bm{U}_j$:
\begin{equation}
    {\hat{y}}_{ij}^{ec}=\sigma(MLP{(\bm{Z}_j}^e||{\bm{Z}_i}^c||RPE_{ij}).
\end{equation}
${\hat{y}}_{ij}^{ec}$ represent the binary classification logits indicating whether $\bm{U}_i$ is the cause of $\bm{U}_j$.
Based on ${\hat{y}}_{ij}^{ec}$, we can determine whether $\bm{U}_j$, $\bm{U}_i$ and $\hat{y}_j^e$ can form a true utter-cause-emotion triplet.
\renewcommand{\arraystretch}{0.8}
\begin{table*}[htbp]
    \centering
    \scriptsize
    \begin{tabular}{llllccccccccc}
        \toprule
        \textbf{Dataset} & \multicolumn{2}{c}{\textbf{Method}} & \textbf{Modality} & \textbf{Anger} & \textbf{Disgust} & \textbf{Fear} & \textbf{Joy} & \textbf{Sadness} & \textbf{Surprise} & \textbf{6 Avg.} & \textbf{4 Avg.} \\
        \midrule
        \multirow{13}{*}{ECF} 
        & Pipline & MC-ECPE-2steps\textsuperscript{$\triangle$} & T, A, V & 24.39 & 0.00 & 0.71 & 38.84 & 21.60 & 40.24 & 29.32 & 31.92 \\
        \cmidrule{2-12}
        & \multirow{5}{*}{E2E} 
        & ECPE-2D\textsuperscript{$\triangle$} & T & 25.13 & 0.00 & 0.00 & 41.25 & 21.62 & 43.24 & 30.80 & 33.55 \\
        & & RankCP & T & 28.29 & 12.03 & 3.52 & 38.69 & 22.17 & 37.67 & 30.58 & 32.48 \\
        & & UECA-Prompt\textsuperscript{$\triangle$} & T &  27.37 & 12.85 & 7.91 & 37.96 & 22.51 & 39.53 & 30.75 & 32.49 \\
        & & SHARK\textsuperscript{*} & T & 28.65 & 10.42 & 5.33 & 40.41 & 25.35 & 40.45 & 32.24 & 34.33 \\
        & & HiLo\textsuperscript{*}& T, A, V &-&-&-&-&-&-& 33.04 & 35.81 \\
        \cmidrule{2-12}
        & \multirow{1}{*}{LLMs}
        & GPT-4o (5-shots) & T & 28.49 & 17.76 & 12.35 & 31.11 & 27.27 & 33.89 & 29.13 & 30.30\\
       \cmidrule{2-12}
       & \multirow{4}{*}{}
       & \multirow{4}{*}{M\textsuperscript{3}HG (ours)} 
       & T & 34.47 & 18.17 & 12.72 & 43.28 & \underline{32.22} & 45.82 &37.46 & 39.95 \\
       & & & T, A & \underline{35.53} & \underline{18.71}  & \underline{17.07}  & \underline{47.73} & 30.97 & 46.72 & \underline{39.10} & \underline{40.97} \\
       & & & T, V & 34.05 & 18.18 & \textbf{19.57} & 46.23 & 32.10 & \textbf{48.50} & 38.90 & 40.72 \\
       & & & T, A, V & \textbf{36.08} & \textbf{23.33} & 9.88 & \textbf{49.03} & \textbf{32.41}  & \underline{47.46} & \textbf{40.07} & \textbf{41.96} \\

        \cmidrule{1-12}
        \multirow{13}{*}{MECAD} 
        & Pipeline & MC-ECPE-2steps & T, A, V & 28.43 & 0.00 & 0.23 & 22.45 & 27.67 & 45.14 & 22.01 & 24.83 \\
        \cmidrule{2-12}
        & \multirow{5}{*}{E2E} 
        & ECPE-2D & T & 28.12 & 0.00 & 0.56 & 24.30 & 28.01 & 35.87 & 25.32 & 28.54 \\
        & & RankCP & T & 29.79 & 12.50 & 3.06  & 21.79 & 29.31 & 32.36 & 26.29 & 28.32 \\
        & & UECA-Prompt & T &  28.54 & 12.12 & 5.32 & 20.84 & 29.67 & 34.17 & 25.91 & 27.87 \\
        & & SHARK & T & 30.22 & 10.16 & 4.10 & 25.84 & 30.21 & 34.59 & 27.58 & 29.99 \\
        & & HiLo\textsuperscript{*} & T, A, V &-&-&-&-&-&-& - & - \\
        \cmidrule{2-12}
        & \multirow{1}{*}{LLMs}
        & GPT-4o (5-shots) & T & 36.65 & 20.08 & 8.45 & 24.52 & 17.89 & 39.77 & 27.16 & 28.42 \\
       \cmidrule{2-12}
       & \multirow{4}{*}{}
       & \multirow{4}{*}{M\textsuperscript{3}HG (ours)} 
       & T & 35.85 & 18.05 & 15.38 & 25.95 & 29.13 & 42.11 & 30.81 & 32.55\\
       & & & T, A & \underline{37.29} & \underline{21.03} & \underline{15.89} & \underline{27.15}  &  30.34 & 42.78 & \underline{32.16} & \underline{33.73} \\
       & & & T, V & 36.91 & 20.48 & \textbf{16.91} & 25.47  & \underline{30.96} & \underline{43.14} & 31.95 & 33.52\\
       & & & T, A, V & \textbf{38.34} & \textbf{21.89} & 8.79 & \textbf{28.10} &  \textbf{31.17} & \textbf{43.29} & \textbf{32.82} & \textbf{34.59} \\
        \bottomrule
    \end{tabular}
    \caption{Performance comparison of different methods on the MECTEC task.
    $\triangle$ denotes the results are from~\cite{wang2023generative}. $*$ denotes the results are from the original paper~\cite{wang2023generative, li2024multimodal}. The best results and the second best results are in bold and underlined, respectively. Since HiLo~\cite{li2024multimodal} is not publicly available, we only report the results of HiLo on the ECF dataset.}
    \label{table:exp}
    \vspace{-2mm}
\end{table*} 

\subsection{Training}
We use Focal loss~\cite{ross2017focal} to cope with category imbalance in emotion-cause classification.
Specifically, the loss of both emotion prediction and cause utterance prediction and the emotion-cause pair prediction, can be expressed as:
% \begin{small}
\begin{equation}
    \mathcal{L}^\beta = -\frac{1}{N^\beta} \sum_{i=1}^{N^\beta} \alpha^\beta (1 - \hat{y}_i^\beta)^\gamma \log(\hat{y}_i^\beta),\\
    \beta \in \{e, c, ec\}
\end{equation}
% \end{small}
where $\beta$ represents the task type, $N^\beta$ denotes the corresponding sample number of $\beta$, $\alpha^\beta$ is the category balancing factor, and $\gamma$ denotes the Focal loss modulation parameter.
These three training losses are optimized jointly during the training process.
\section{EXPERIMENTS}
\subsection{Experimental Settings}
We conduct extensive experiments on two MECTEC benchmark datasets, i.e., \textbf{ECF}~\cite{ecf} and \textbf{MECAD}, which both contain data of three modalities: text, audio, and video.
Similar to~\cite{ecf}, we evaluate the model's overall performance using the F1 score.
The F1 score is computed for utter-cause-emotion triplets within each emotion category separately.
Then the weighted average F1 score is calculated across all six emotion categories which is referred to \textit{6 Avg}.
In addition, as in~\cite{wang2023generative}, considering the data imbalance among different emotion categories, we also report the weighted average F1 scores for the four main emotion categories except \textit{Disgust} and \textit{Fear}, which is referred to \textit{4 Avg}. 
The implementation details of the experiment are given in Appendix~\ref{appendix-d}.

\subsection{Baselines}
Due to the limited research on the MECTEC task, representative approaches in related fields of Emotion Cause Pair Extraction (ECPE) and Emotion Cause Pair Extraction in Conversations (ECPEC) are considered.
The ECPE and ECPEC tasks aim to extract emotion-cause pairs from plain texts and conversations, respectively.

We compare our model with seven baselines:
1) \textbf{MC-ECPE-2steps}~\cite{ecf} is a two-step MECTEC architecture, which first extracts emotion utterances and cause utterances separately, and then performs pairing and filtering to identify emotion-cause pairs.
2) \textbf{HiLo}~\cite{li2024multimodal} is one of the SOTA approaches for the MECTEC task, which fully utilizes conversion information through a labeling constraint mechanism.
3) \textbf{ECPE-2D}~\cite{ECPE-2D} is an E2E framework for ECPE that uses 2D-Transformer to model the interactions of emotion-cause pairs.
4) \textbf{RankCP}~\cite{wei2020effective} is a GAT-based approach for ECPE to extract emotion-cause pairs by ranking.
5) \textbf{UECA-Prompt}~\cite{UECA-Prompt} is one of the SOTA methods for ECPE, which decomposes the task into multiple objectives and converts them into sub-prompts.
6) \textbf{SHARK}~\cite{wang2023generative} is the SOTA method for ECPEC that incorporates commonsense into GATs to improve the model's semantic understanding of emotions and causes.
7) \textbf{GPT-4o} is one of the most powerful large language models (LLMs) for open-domain conversations.
% used in a few-shot learning setting for extracting utter-cause-emotion triplets.
Details of prompts are provided in Appendix~\ref{gpt}.

\subsection{Experimental Results}
Table~\ref{table:exp} shows the experimental results of M\textsuperscript{3}HG and seven baseline models evaluated on the ECF dataset and the MECAD dataset.
Our model demonstrates an excellent performance both on the ECF dataset and the MECAD dataset.

\noindent\textbf{Results on the ECF dataset.} First of all, as shown in Table~\ref{table:exp}, among all the baseline models, the E2E approaches such as SHARK and HiLo deliver the best performance, indicating that the E2E framework is more effective compared to the two-step pipeline frameworks.
In contrast, M\textsuperscript{3}HG adopting three modalities outperforms the SOTA E2E model HiLo, with 21.28\% and 17.17\% improvement in \textit{6 Avg} and \textit{4 Avg} scores, respectively.
We attribute this improvement to M\textsuperscript{3}HG’s ability to effectively extract semantic information at inter-utterance and intra-utterance levels, which enables the model to accurately pair emotion utterances and cause utterances.
Specifically, in two challenging emotion categories which have limited training samples, i.e. \textit{Disgust} and \textit{Fear}, M\textsuperscript{3}HG also exhibits high performances.
For example, compared to GPT-4o, which achieved the second highest F1 scores in the \textit{Disgust} and \textit{Fear} categories, M\textsuperscript{3}HG shows improvements of 31.36\% and 58.46\%, respectively.

When only incorporating the text modality, the \textit{6 Avg} and the \textit{4 Avg} scores of M\textsuperscript{3}HG are 37.46 and 39.95.
When incorporating audio and video with the text modality separately, the performance of M\textsuperscript{3}HG is improved to 39.10, 38.90 of \textit{6 Avg} scores and 40.97, 40.72 of \textit{4 Avg} scores.
When incorporating all three modalities, M\textsuperscript{3}HG achieves the highest performance with 40.07 of \textit{6 Avg} scores and 41.96 of \textit{4 Avg} scores.
Meanwhile, it can be observed that M\textsuperscript{3}HG outperforms all the baseline models even when only using the text modality, demonstrating its superiority on the ECF dataset.

\noindent\textbf{Results on the MECAD dataset.} As shown in Table~\ref{table:exp}, M\textsuperscript{3}HG also achieves the highest results on the MECAD dataset.
Compared to the second best model SHARK, M\textsuperscript{3}HG adopting three modalities achieves the improvement of 19\% on the \textit{6 Avg} scores and 15.34\% on the \textit{4 Avg} score.
Furthermore, despite GPT-4o’s superior semantic comprehension abilities, its performance on the MECAD dataset remains suboptimal, with its \textit{6 Avg} score and \textit{4 Avg} score of 27.16 and 28.42.
Therefore, the few-shot-based LLM approach still struggles to effectively handle the MECTEC task.
As shown in Table~\ref{table:exp}, M\textsuperscript{3}HG exhibits a universal highest performance on the MECAD dataset, demonstrating the superiority and robustness of M\textsuperscript{3}HG when dealing with multi-conversation scenarios.

% 有关M3HG在不同对话场景下更详细的实验结果见附录
More detailed experimental results and the ablation study on M\textsuperscript{3}HG are presented in Appendix~\ref{more experimental results}.


\section{Conclusion}
In this work, we propose the first multimodal and multi-scenario Chinese emotion-cause analysis dataset, MECAD, for MECTEC and related emotion cause analysis tasks.
Compared to ECF, the only existing dataset for multimodal emotion-cause analysis, MECAD offers more diverse conversation scenarios.
It helps to enhance the generalizability and applicability of MECTEC models in complex social media environments.
Moreover, MECAD is a valuable resource for cross-cultural emotion analysis and recognition.
Furthermore, we propose a generalized MECTEC framework named M\textsuperscript{3}HG, which deeply extracts emotional and causal contexts, while effectively integrating semantic information across multiple granular levels.
Extensive experiments on the ECF dataset and the MECAD dataset demonstrate the superiority of our method compared to the existing state-of-the-art methods.
\section*{Limitations}
There are also some potential limitations in this work.
% First, the structural complexity of the M\textsuperscript{3}HG heterogeneous graph depends on the predefined local context window $K$.
% With different $K$, The model’s performance shows significant fluctuation.
% The determination of $K$ requires more  observations from further experiments.
% 首先，M3HG无法处理过长的对话，其可输入的最大对话长度受到使用的语言模型限制。
First, the process of emotional and causal context extraction does not integrate external knowledge, which limits the model's accuracy for emotion prediction and cause prediction.
In the future, we plan to integrate external knowledge into our model and leverage the advanced semantic extraction capabilities of current LLM technology to facilitate deeper and more precise emotion cause analysis. 
Second, M\textsuperscript{3}HG cannot handle excessively long conversations, as its input length is constrained by the language model used.
Furthermore, M\textsuperscript{3}HG may suffer from error propagation in the multimodal fusion process when emotion labels have uneven information across modalities.
This imbalance can lead to inaccurate predictions, especially when modalities conflict. This challenge is common in current multimodal models for emotion-cause analysis and suggests an area for future improvement. 

\section*{Ethical Considerations}
% 我们在收集数据时没有采用真实世界的对话，这是因为这类对话可能侵犯到说话人的隐私。
We did not use real-world conversations in our data collection because such conversations may violate the privacy of the speaker.
% 招募演员进行扮演的效果和TV series相同，但场景不如TV series多样。
The effect of recruiting actors to play the roles is the same as in the TV series, but the scenes are not as diverse as in the TV series.
% 因此，我们使用TV series作为数据源。
Therefore, we use TV series as the data source.
% 除此之外，在进行数据标注时我们也对标注内容进行了匿名化和去标识化，有效保证了数据收集过程符合伦理标准。
To further protect privacy, all data annotations were anonymized and de-identified, ensuring that our data collection adheres to ethical standards.
% Bibliography entries for the entire Anthology, followed by custom entries
%\bibliography{anthology,custom}
% Custom bibliography entries only
\bibliography{custom}

\appendix


\section{Dataset Statistics and Analysis of MECAD \label{appendix-a}}




% 为了保证MECAD的标注质量，我们计算了每两个标注者间共同标注的数据的Cohen's kappa分数，如图所示。
To ensure the annotation quality of MECAD, we calculated Cohen's kappa~\cite{cohens-kappa} scores for every co-annotated data between two annotators, as shown in Figure~\ref{fig:annotation result}. 
% 我们可以观察到，所有标注者之间的Cohen's kappa~\cite{cohens-kappa} scores均在0.6左右，显示出良好的标注一致性。
The Cohen’s kappa~\cite{cohens-kappa} scores across all annotators are consistently around 0.6, indicating a good level of annotation consistency.

After the labeling was completed, we computed Cohen's kappa scores separately for data that were not co-labeled between the two labelers, as shown in Figure~\ref{fig:annotation result}.
Table~\ref{tab:MECAD statistics} lists some statistics of the MECAD dataset.
The dataset contains a total of 989 conversations, 10,516 utterances, and 8,077 emotion cause pairs from 56 different TV series, which ensures the size and diversity of the dataset.
Similar to M\textsuperscript{3}ED~\cite{M3ED}, we used TV-independent data segmentation to ensure the ability to validate model robustness as a benchmark dataset.
The average number of utterances and the average length of an utterance of a conversation are similar in the training, validation, and test sets.
At the same time, we can find that the average relative positions of the emotion cause pairs are all around 1, indicating that most of the emotions in the conversation are caused by the previous utterance.

\begin{figure}
\centering
  \includegraphics[width=0.45\textwidth]{Figure/annotationtool.png}
  \caption{The interface of the developed online multimodal conversation emotion cause annotation toolkit.}
  \label{fig:annotation toolkit}
\end{figure}


\begin{figure}
\centering
  \includegraphics[width=0.45\textwidth]{Figure/annotators_kappa_matrix.pdf}
  \caption{Schematic representation of Cohen's Kappa scores for the common labeled portion between every two annotators. A blank section indicates that there is no common annotation data between two annotators.}
  \label{fig:annotation result}
\end{figure}
\begin{table}
  \caption{MECAD statistics. \textit{Rel pos of ec pairs} denotes the relative position between emotion utterances and cause utterances in emotion-cause pairs.}
  \label{tab:MECAD statistics}
  \scriptsize
  \centering
  \begin{tabular}{lccc|c}
    \toprule
    Statistic & Train & Val & Test & Total\\
    \midrule
    \# TV series & 38 & 7 & 11 & 56\\
    \# conversations & 684 & 126  & 179 &  989 \\ 
    \# uttrs & 7,516 & 1,168  & 1,832 & 10,516\\
    \# spkrs & 421 & 87 & 118 & 626 \\
    Avg. uttrs/conversation & 10.99 & 9.27 & 10.24 & 10.63 \\
    Avg. uttr length & 18.30 & 18.80 & 18.15 & 18.33\\
    Avg. rel pos of ec pairs & 0.72 & 0.73 & 0.55 & 0.69\\
    Max. rel pos of ec pairs & 13 & 7 & 6 & 13\\
    Min. rel pos of ec pairs & -14& -5  & -9 & -14\\
    Emotion uttrs with cause & 4,526 & 743  & 1,062 & 6,331 \\
    ec pairs & 5,788  & 977  & 1,312 & 8,077 \\
  \bottomrule
  \end{tabular}
\end{table}
\label{appendix-b}
\begin{figure}
\centering
  \includegraphics[width=0.5\textwidth]{Figure/cause_type_distribution.pdf}
  \caption{Percentage of five cause types in the MECAD dataset and percentage of modal basis for emotion cause inferences.}
  \label{fig:cause_type_distribution}
\end{figure}
% 我们参考ECF[14]对情感原因的分类方式，并进一步增加了一个被称为Self Reflection的类别，该类别与其余四个类别不同，表示情绪可能由个人的内省或自我反思引发，如对过去事件的回忆或对未来的担忧。
We referred how the ECF~\cite{ecf} dataset categorizes the emotion causes and added a new category called \textit{Self Reflection}, which differs from the remaining four categories by indicating that emotions may be triggered by an individual's introspection or self-reflection, such as recollections of past events or worries about the future.
% 如图三所示，event类型为占比最大的原因类型，表示大多数情绪都由对话中具体的事件引起。
As shown in Figure~\ref{fig:cause_type_distribution}, the event type is the cause type with the largest share, indicating that most of the emotions are caused by specific events in the conversation.
% 值得注意的是，在我们的数据集中，有36.1%的情绪原因反映在音频和视频模态中，这充分体现了多模态场景研究的必要性。
Notably, 36.1\% of the causes of emotion in our dataset are reflected in both audio and video modalities, which exemplifies the need for multimodal scene studies.

\section{The Annotation Toolkit of MECAD}
\label{Appendix-toolkit}
To enhance annotation efficiency and accuracy, we developed an online multimodal conversation emotion cause annotation tool based on web technology\footnote{The annotation tool has been open-sourced at \url{https://github.com/redifinition/MECAD-MECTEC}}.
% 如图所示，该工具包的主页展示了不同用户需要标注的对话列表以及标注进度。
As illustrated in Figure~\ref{fig:annotation toolkit}, the toolkit’s homepage presents a list of conversations assigned to the corresponding annotators, along with the progress of their annotations.
The conversation annotation page displays speaker information, video segments, corresponding scripts, and configurable annotation items, enabling annotators to quickly and efficiently complete their annotations.

% 通过灵活可修改的页面，相关研究人员可以利用我们的标注工具用于后续其他多模态情感分析研究的数据集构建中。
With flexible and modifiable web pages, researchers can utilize our annotation tools in dataset constructions for further multimodal sentiment analysis studies.
% % 这是一个用于对话多模态标注的通用工具包，具有标注内容快速配置、多用户与任务管理、标注一致性管理、可分布式部署等功能。图2展示了标注工具的界面。
% This is a generalized toolkit for multimodal annotation of conversations, with features such as quick configuration of annotation content, multi-user and task management, annotation consistency management, and distributable deployment. 
\begin{algorithm}
\caption{Super-Node-based Graph Constructuon for a Conversation}
\footnotesize
\begin{algorithmic}[1]
\State \textbf{Input}: the conversation $\{S_1:U_1, S_2:U_2, ..., S_N:U_N\}$, speaker identity $p(\cdot)$ satisfies $p(U_i) = S_i$, the direct context window $K$
\State \textbf{Output}: Super-Node-based M$^3$HG: $\mathcal{G} = (\mathcal{V}, \mathcal{E}, \mathcal{R})$
\State $\mathcal{V} \gets \{(SN_u^1, N_e^1, N_c^1), ..., (SN_u^N, N_e^N, N_c^N), SN_d^1\}$
\State $\mathcal{E} \gets \emptyset$
\State $\mathcal{R} \gets \{r_{ss}, r_{ds}, r_{gc}, r_{ec}, r_{cc}\}$
\For{$i \in \{2, 3, ..., N\}$}
    \State $c \gets 0$, $w \gets i - 1$
    \While{$w > 0$ and $c < K$} 
        \If{$p(U_w) = p(U_i)$}
            \State $\mathcal{E} \gets \mathcal{E} \cup \{(SN_u^w, SN_u^i, r_{ss})\}$ 
            \State $c \gets c + 1$
        \Else
            \State $\mathcal{E} \gets \mathcal{E} \cup \{(SN_u^w, SN_u^i, r_{ds})\}$ 
        \EndIf
        \State $w \gets w - 1$
    \EndWhile
\EndFor
\For{$i \in \{1, 2, ..., N\}$}
    \State $\mathcal{E} \gets \mathcal{E} \cup \{(SN_u^i, N_e^i, r_{ec})\}$ 
    \State $\mathcal{E} \gets \mathcal{E} \cup \{(SN_u^i, N_c^i, r_{cc})\}$ 
    \State $\mathcal{E} \gets \mathcal{E} \cup \{(SN_u^i, SN_d^i, r_{gc})\}$ 
\EndFor
\State \textbf{return} $\mathcal{G} = (\mathcal{V}, \mathcal{E}, \mathcal{R})$
\end{algorithmic}
\end{algorithm}
\section{\texorpdfstring{Design Details of M\textsuperscript{3}HG}{Design Details of M3HG}}
\subsection{Multimodal Feature Extracting}
\label{appendix-c.1}
% Text : 我们将对话中的所有文本模态的话语及对应的说话人进行拼接，并且加入若干特殊令牌，得到文本模态的输入序列：
\textbf{Text} : We splice all the textual modal utterances and the corresponding speakers in the conversation and add a number of special tokens to get the textual modal input sequence:
$X^t = \{<cls\_token> S_1:U_1^t, <sep\_token>, \dots, <cls\_token>S_n:U_n^t, <sep\_token> \}$,
% 其中<cls_token>和<sep_token>表示预训练语言模型（PLM）的起始和结束令牌。
where $<cls\_token>$ and $<sep\_token>$ denote the classification token and the separation token used in the pre-trained language model (PLM), respectively.
% 为了让超出PLM最大输入序列长度的对话在输入PLM时尽可能的保留上下文信息，我们在预处理时依次截断对话中最大长度的话语的最后一个标记，直到满足PLM的最大序列长度要求为止。
To allow conversations that exceed the maximum input sequence length of the PLM to retain as much contextual information as possible when they are fed into the PLM, we sequentially truncate the last tokens of the maximum-length utterances of the conversation during preprocessing until the maximum sequence length requirement of the PLM is met.
% 然后将输入序列X^t输入到PLM中以获得整个对话的序列表示：
The input sequence $X^t$ is then fed into the PLM to obtain a sequential representation of the entire conversation:
\begin{equation}
\bm{I}^t = PLM(X^t),
\end{equation}

% 其中 $I^t \in \mathbb{R}^{L \times d_t}$，其中 $L$ 是输入序列的长度，$d_t$ 是 PLM 的隐层维度。
where $\bm{I}^t \in \mathbb{R}^{L \times d_t}$, $L$ is the length of the input sequence and $d_t$ is the hidden dimension of the PLM.
% 为了获得每个话语的序列表示，我们将I^t中每个对话的标记的序列表示作加权平均，得到每个话语的序列表示E^t∈R^(N×d_t ),其中N表示该对话的话语数。
To obtain the sequence representation of each utterance, we make a weighted average of the sequence representations of the tokens of each conversation in $\bm{I}^t$ to obtain the sequence representation of each utterance $\bm{E}^t \in \mathbb{R}^{N \times d_t}$,where $N$ denotes the number of utterances of that conversation.
% 我们选取Speaker-Aware RoBERTa (SA-RoBERTa)[49]作为PLM。
We selected Speaker-Aware RoBERTa (SA-RoBERTa)~\cite{SA-RoBERTa} as the PLM.

% Audio : 在将音频采样到16khz后，我们其输入到音频特征提取模型(AFE)中得到对话音频模态的序列表示：
\textbf{Audio} : After resampling the audio to 16khz, we input it into an audio feature extraction model (AFE) to get a sequential representation of the audio modality of the conversation:
\begin{equation}
\bm{E}^a = AFE(X^a),
\end{equation}
% 其中E^a∈R^(N×d_a )，d_a是音频特征提取模型的隐层维度。我们选取Wav2Vec2.0[50]作为音频特征提取模型。
where $\bm{E}^a \in \mathbb{R}^{n \times d_a}$, and $d_a$ is the hidden layer dimension of the audio feature extraction model. 
We choose Wav2Vec2.0~\cite{wav2vec2} as the audio feature extraction model.

% Video: 我们首先将视频进行等间隔采样为若干图片序列，得到视频模态的输入序列X^v, X^v∈R^(F×d_f×d_f )，其中F为采样后的帧数，d_f为图片的尺寸。
\textbf{Video}: We first sample the video at equal intervals as a sequence of images over several frames to obtain the input sequence $\bm{X}^v$, $\bm{X}^v \in \mathbb{R}^{F \times d_f \times d_f }$ of the video modality, where $F$ is the number of sampled frames and $d_f$ is the size of the picture.
% 然后将图片序列输入到视频特征提取模型(VFE)中得到视频模态的序列表示：
The image sequences are then fed into the video feature extraction model (VFE) to get a sequence representation of the video modalities:
\begin{equation}
\bm{E}^v = VFE(\bm{X}^v),
\end{equation}
where $\bm{E}^v \in \mathbb{R}^{n \times d_v}$ and $d_v$ is the hidden layer dimension of the video feature extraction model.
We select the pre-trained DenseNet\\~\cite{DenseNet} as the audio feature extraction model.

\subsection{Pseudo-code of Graph Construction}
\label{appendix:pseudo-code}

The pseudo-code of the graph construction process is shown in Algorithm~\ref{appendix:pseudo-code}.

\begin{table*}
  \centering
  \caption{Performance comparison of different methods for conversations with varying numbers of utterances. The best results and the second best results are in bold and underlined, respectively.}
  \label{tab:conversation length}
  \renewcommand{\arraystretch}{1.05}
  \small  % 设置表格与正文一致的字体大小
  {%
  \begin{tabular}{lccccccccc}
  \toprule
  \multirow{4}{*}{Method} & \multicolumn{4}{c}{\textbf{ECF}} & \multicolumn{4}{c}{\textbf{MECAD}} \\ \cmidrule(l){2-5} \cmidrule(l){6-9} & \multicolumn{2}{c}{num\_utt $\le$ 10} & \multicolumn{2}{c}{num\_utt $>$  10} & \multicolumn{2}{c}{num\_utt $\le$ 10} & \multicolumn{2}{c}{num\_utt $>$ 10} \\ \cmidrule(l){2-3} \cmidrule(l){4-5} \cmidrule(l){6-7} \cmidrule(l){8-9}
   & 6 Avg. & 4 Avg. & 6 Avg. & 4 Avg. & 6 Avg. & 4 Avg. & 6 Avg. & 4 Avg. \\ 
   \midrule
  RankCP  &  31.50 & 33.29 & 29.34 & 31.88 & 27.19 & 29.23 & 25.11 & 27.13 \\
  SHARK & 33.68 & 35.57 & 31.49 & 33.17 &  28.32 &   30.75 & 27.01 & 29.41 \\
  GPT-4o & 30.08 & 31.56 & 28.42 & 29.36 & 26.34 & 27.82 & 27.79 & 28.88 \\
  M\textsuperscript{3}HG (T) & \underline{39.18} & \underline{41.25} & \underline{36.18} & \underline{38.98} & \underline{31.95} & \underline{33.40} & \underline{29.94} & \underline{31.90} \\
M\textsuperscript{3}HG (T, A, V) & \textbf{41.95} & \textbf{40.42} & \textbf{38.67} & \textbf{41.09} & \textbf{33.76} & \textbf{35.21} & \textbf{32.10} & \textbf{34.12} \\
  \bottomrule
  \end{tabular}%
  }
\end{table*}
\subsection{An Example of the Graph construction}
\label{appendix-c.3}
\begin{figure}
  \includegraphics[width=0.9\linewidth]{Figure/GraphExample.pdf}
  \caption{Super-Node-based edges and relations constructed from a conversation in MECAD with K = 1. The utterance Super-Nodes of the two speakers are shown in gray and blue, respectively. The black solid and dashed lines denote the Super-Edges between the same speaker and different speakers, respectively, and the red dotted lines denote the Super-Edges between the utterance Super-Nodes and the conversation Super-Nodes.}
  \label{fig:graph-example}
\end{figure}
% 若取K=1，对于图一中的对话可构建如图四所示的图结构。
If $K=1$, the graph constructed for the conversation in Figure~\ref{fig:introduction figure} is shown in Figure~\ref{fig:graph-example}.


\section{Implement Details of the Experiment}
For the ECF dataset, we use the pre-trained RoBERTa-large\footnote{\url{https://huggingface.co/FacebookAI/roberta-large}} model to initialize the feature extraction parameters of the text modality.
For audio modality, we use the wav2vec2-base-960h\footnote{\url{https://huggingface.co/facebook/wav2vec2-base-960h}} model and for video modality we use the DenseNet~\cite{DenseNet} model.
For the MECAD dataset, we use the chinese-roberta-wwm-ext-large\footnote{\url{https://huggingface.co/hfl/chinese-roberta-wwm-ext-large}} model for the initialization of textual modal features, the wav2vec2-large-chinese-zh-cn\footnote{\url{https://huggingface.co/wbbbbb/wav2vec2-large-chinese-zh-cn}} model for the extraction of audio modal features, and the DenseNet model is also applied to the video modal.
In our experiments, none of the parameters of the PLM were frozen.
During the construction of the graph, we set the hyperparameter $K$ to 3.
During training, we use the AdamW~\cite{AdamW} optimizer with batch size and learning rate set to 16 and 5e-6, respectively, and perform a parameter update after every two mini-batches.
Our model is trained for 50 epochs on the training set, and the checkpoints corresponding to the highest values of the weighted average F1 scores of the six emotions on the validation set are used as the results of the test set.
\label{appendix-d}

\begin{table*}
  \centering
  \caption{Performance comparison of different methods on four subtasks. The best results and the second best
results are in bold and underlined, respectively.}
  \label{tab:sub_task}
  \renewcommand{\arraystretch}{1.05}
  \scriptsize
  \begin{tabular}{lllcccccccccccc}
  \toprule
  \multirow{2}{*}{Dataset} & \multirow{2}{*}{Method} & \multicolumn{3}{c}{EP} & \multicolumn{3}{c}{ER} & \multicolumn{3}{c}{CE} & \multicolumn{3}{c}{EC} \\
  \cmidrule(lr){3-5} \cmidrule(lr){6-8} \cmidrule(lr){9-11} \cmidrule(lr){12-14}
   & & P & R & F1 & P & R & F1 & P & R & F1 & P & R & F1 \\
  \midrule
  
  \multirow{3}{*}{ECF} 
    & SHARK             & \underline{59.00} & 61.21 & \underline{60.74} & \underline{40.34} & \textbf{45.65} & \underline{42.83} & \underline{69.25} & 66.13 & 67.64 & \underline{50.12} & 46.31 & \underline{48.14} \\
    & GPT-4o(5-shots)   & 45.17 & \textbf{78.62} & 57.37 & 36.42 & 42.70 & 36.76 & 57.02 & \textbf{84.85} & \underline{68.21} & 32.90 & \textbf{61.13} & 42.78 \\
    & M\textsuperscript{3}HG (T)     & \textbf{71.36} & \underline{75.11} & \textbf{73.19} & \textbf{52.24} & \underline{45.63} & \textbf{46.60} & \textbf{72.32} & \underline{68.40} & \textbf{70.30} & \textbf{58.03} & \underline{52.05} & \textbf{54.88} \\
  
  \midrule
  
  \multirow{3}{*}{MECAD}
    & SHARK             & 69.30 & \underline{67.02} & \underline{68.14} & \underline{39.38} & \underline{36.93} & \underline{38.12} & 64.18 & 66.36 & 65.24 & \underline{49.02} & \underline{42.87} & \underline{45.74} \\
    & GPT-4o(5-shots)   & \underline{71.41} & 63.69 & 67.33 & 38.69 & 36.59 & 34.22 & \underline{65.03} & \underline{69.26} & \underline{67.08} & 39.68 & 41.77 & 40.70 \\
    & M\textsuperscript{3}HG (T)     & \textbf{72.34} & \textbf{67.84} & \textbf{70.02} & \textbf{43.35} & \textbf{40.98} & \textbf{41.66} & \textbf{66.12} & \textbf{70.24} & \textbf{68.12} & \textbf{54.82} & \textbf{46.42} & \textbf{50.27} \\
  
  \bottomrule
  \end{tabular}
\end{table*}

\section{\texorpdfstring{Supplementary Experimental Results of M\textsuperscript{3}HG}{Supplementary Experimental Results of M3HG}}
\label{more experimental results}
\subsection{Ablation Study}
\label{ablation study}
\textbf{Effect of different modules.} We conduct ablation studies to verify the effectiveness of different modules in M\textsuperscript{3}HG on the two datasets using \textit{6 Avg} and \textit{4 Avg} scores.
As shown in Table~\ref{tab:ablation study}, $w/o$ $N^e \& N^c$ indicates no use of emotional and causal context nodes in graph construction.
Consequently, emotion-cause pair prediction is performed directly based on the features of each utterance node.
$w/o$ \textit{inter-fusion} and $w/o$ \textit{intra-fusion} denote the absence of inter-utterance and intra-utterance multimodal fusion, respectively, during multi-scale semantic information fusion.
Our model outperforms the state-of-the-art baselines even without utilizing the previous three mechanisms.
Specifically, the performance of M\textsuperscript{3}HG degrades on both ECF and MECAD datasets when removing the emotional and causal context nodes, demonstrating the necessity of explicitly modeling the emotion and cause-related contexts.
Moreover, removing both intra-utterance and inter-utterance semantic fusion results in a drop in the model's performance, the former of which causes a more significant degradation.
It highlights the importance of effectively fusing semantic information at different scales within heterogeneous graphs, particularly within individual utterances.

\begin{table}
  \caption{Ablation results.}
  \scriptsize
  \label{tab:ablation study}
  \centering
  \begin{tabular}{clll}
    \toprule
    Dataset & Model & 6 Avg. &  4 Avg.\\
    \midrule
    \multirow{5}{*}{ECF} & M\textsuperscript{3}HG & \textbf{40.07} & \textbf{41.96} \\
    & $w/o$ all modules & $36.81_{(\downarrow3.26)}$ & $38.57_{(\downarrow3.39)}$ \\
    & \quad $w/o$ $N^e \& N^c$ & $38.13_{(\downarrow1.94)}$& $40.11_{(\downarrow1.85)}$\\
    & \quad $w/o$ \textit{inter-fusion} & $39.56_{(\downarrow0.51)}$ & $41.14_{(\downarrow0.82)}$ \\ 
    & \quad $w/o$ \textit{intra-fusion} &$39.12_{(\downarrow0.95)}$ & $40.86_{(\downarrow1.10)}$ \\ 
    \midrule
    \multirow{5}{*}{MECAD} & M\textsuperscript{3}HG & \textbf{32.82} & \textbf{34.59} \\
    & $w/o$ all modules & $30.37_{(\downarrow2.45)}$ & $32.27_{(\downarrow2.32)}$\\
    & \quad $w/o$ $N^e \& N^c$ & $30.94_{(\downarrow1.88)}$& $32.79_{(\downarrow1.80)}$\\
    & \quad  $w/o$ \textit{inter-fusion} & $32.57_{(\downarrow0.25)}$ & $33.91_{(\downarrow0.68)}$ \\ 
    & \quad $w/o$ \textit{intra-fusion} & $32.16_{(\downarrow0.66)}$ & $33.33_{(\downarrow1.26)}$ \\ 
  \bottomrule
  \end{tabular}
\end{table}

\begin{figure}
  \centering
  \includegraphics[width=0.7\linewidth]{Figure/K.pdf}
  \caption{Results of M\textsuperscript{3}HG with various $K$ values.}
  \label{fig:k-value}
\end{figure}

\noindent \textbf{Effect of the hyperparameter $K$.}
% 超参数K的大小与M3HG图结构构建的时空复杂度密切相关。
The hyperparameter $K$ is closely related to the spatio-temporal complexity of the M\textsuperscript{3}HG's graph construction.
% 我们改变K的大小（从 1 到 6 不等）来测试其效果，模型在两个数据集的性能如图所示。
We vary the size of K (ranging from 1 to 6) to test its effect, and the result of the M\textsuperscript{3}HG on both datasets is shown in Figure~\ref{fig:k-value}. 
The performance of M\textsuperscript{3}HG on both datasets initially improves with increasing $K$ and then declines, with the best performance observed at $K = 3$.

\begin{figure*}
  \centering
  \includegraphics[width=\linewidth]{latex/Figure/success_case.pdf}
  \caption{Comparison of utter-cause-emotion triplet on two test samples.}
  \label{fig:case study}
\end{figure*}

\subsection{In-Depth Analysis}
\label{in-depth analysis}
\noindent \textbf{The impact of conversation length.} 
To evaluate the performance of M\textsuperscript{3}HG in handling longer conversations, we present a comparison of the performance of M\textsuperscript{3}HG and other baseline models across conversations of varying lengths, as shown in Table~\ref{tab:conversation length}. 
We observe that M\textsuperscript{3}HG outperforms all baseline models in scenarios involving conversations with more than 10 utterances, which account for 42.65\% and 43.38\% of all conversations in the ECF and MECAD datasets, respectively.
% 在长对话中，包括GPT-4o在内的基线模型未能有效提取对话的全局上下文信息，从而忽略了一些triplets.
In long conversations, baseline models, including GPT-4o, fail to effectively extract global contextual information, thereby missing a number of triplets.
Our model, through semantic fusion at different scales within multimodal heterogeneous graphs, effectively captures more triplets by extracting contextual information from long conversations.

\noindent \textbf{Model performance when cause utterance appears after emotion utterances.} 
A key challenge in the MECTEC task is when the cause of a speaker’s emotion is revealed later in the conversation, requiring the model to effectively capture and interpret the global context of the conversation.
To further emphasize the superior performance of M\textsuperscript{3}HG in handling cases where the cause utterance appears after the emotion utterance, we identified and filtered all such conversations from the ECF and MECAD datasets. The performance of M\textsuperscript{3}HG, compared with two other representative models, is shown in Table~\ref{tab:cause appears after emotion}.
M\textsuperscript{3}G demonstrates superior performance, while SHARK suffers a greater performance drop compared to M\textsuperscript{3}HG.
Although the performance drop for GPT-4o (5-shots) is less pronounced, its overall performance remains unsatisfactory.

\begin{table}
  \centering
  \renewcommand{\arraystretch}{1.05}
  \caption{Performance comparison of different methods for conversations in which cause utterance appears after emotion utterances. The best results and the second best results are in bold and underlined, respectively.}
  \label{tab:cause appears after emotion}
  \small  % 设置表格与正文一致的字体大小
  {%
  \begin{tabular}{lcccc}
  \toprule
  \multirow{3}{*}{Method} & \multicolumn{2}{c}{\textbf{ECF}} & \multicolumn{2}{c}{\textbf{MECAD}} \\ \cmidrule(l){2-3} \cmidrule(l){4-5} 
   & 6 Avg. & 4 Avg. & 6 Avg. & 4 Avg.  \\ 
   \midrule
  SHARK & 29.15 & 30.54 & 25.49 & 27.51 \\
  GPT-4o & 28.21 & 29.45 & 26.42 & 27.76 \\
  M\textsuperscript{3}HG (T) & \underline{35.48} & \underline{37.01} & \underline{29.18} & \underline{30.93} \\
M\textsuperscript{3}HG (T, A, V) & \textbf{38.25} & \textbf{40.01} & \textbf{31.27} & \textbf{33.09} \\
  \bottomrule
  \end{tabular}%
  }
\end{table}

\noindent \textbf{Model performance on four subtasks.}
To evaluate M\textsuperscript{3}HG’s performance more comprehensively, we define the following four subtasks:
\begin{itemize}[itemsep=0pt,parsep=0pt,topsep=0pt,partopsep=0pt]
  \item \textbf{Emotion Extraction (EP)}: Predict whether an utterance expresses an emotion (binary classification), same as SHARK.
  \item \textbf{Emotion Recognition (ER)}: Predict the emotion category of an utterance (multi-class classification).
  \item \textbf{Cause Extraction (CE)}: Predict whether an utterance is a cause utterance (binary classification), same as SHARK.
  \item \textbf{Emotion-Cause Pair Extraction (EC)}: Predict whether two utterances of a conversation form an emotion-cause pair (binary classification).
\end{itemize}

Table~\ref{tab:sub_task} demonstrates the performance comparison between M\textsuperscript{3}HG and other SOTA models across the four subtasks.
For the EP subtask, M\textsuperscript{3}HG performs the best across both datasets.
It is worth noting that GPT-4o (5-shots) achieves a high recall on the ECF dataset. 
This phenomenon can be attributed to the more pronounced label sparsity in the ECF dataset compared to MECAD. 
As a result, GPT-4o (5-shots) frequently predicts that an utterance carries emotion, leading to a higher recall.
For the ER subtask, M\textsuperscript{3}HG achieves the best results across both datasets.
This demonstrates M3HG’s ability to effectively extract the emotional context embedded in utterances.
For the CE subtask, M\textsuperscript{3}HG performs best, demonstrating the importance of integrating the cause prediction subtask into the model during training.
For the EC subtask, GPT-4o (5-shots) similarly exhibits high recall on the ECF dataset. 
This is due to the severe label sparsity problem in the ECF dataset, compared to MECAD, which leads GPT-4o to predict as many emotion-cause pairs as possible.

\begin{table*}[t]
    \renewcommand{\arraystretch}{1.05}
  \caption{An example of prompt for ChatGPT.}
  \scriptsize
  \label{tab:prompts}
  \centering
  \begin{tabular}{lll}
    \toprule
    \multirow{28}{*}{Input} & \multirow{8}{*}{Instruction} & You are an expert in sentiment analysis and identification of emotional causes. I will give you a conversation \\
    & & between multiple speakers. You are required to extract the utter-cause-emotion triplet for a given utterance.  \\
    & &  First, infer the emotion label for the utterance (select one from: Anger, Disgust, Fear, Joy, Sadness, Surprise \\
    & &   or Neutral). Then, identify the index(es) of the cause utterance(s) that triggered this emotion (the index  \\
    & & should represent the utterance(s) from the conversation that caused the emotion, and it must be non- \\
    & &  negative. Multiple indices  should be separated by commas). If the predicted emotion is Neutral, there is no \\
    & & corresponding cause utterance. The output should follow the format: emotion label, cause utterance indices. \\
    & &  Examples of the expected output format: Example 1: happy,3. Example 2: sad,3,4,5. Example 3: neutral. \\
    \cline{2-3}
    & \multirow{20}{*}{Demonstrations}& Input Conversation :\\
    & & \{ 1. Fang Sijin: First, change your clothes, then head to this address. A decoration company will be coming \\ 
    & &  over shortly. You’ll need to supervise their work and see how you can help. \}\\
    & & \{ 2. Zhu Shanshan: Wait, am I really responsible for this? I don’t know anything about decoration. \} \\
    & & \{ 3. Fang Sijin: You’ve been handing out flyers for two days now. Have you gotten any interested customers? \} \\ 
    & & \{ 4. Zhu Shanshan: But you only told me to distribute the flyers; you never ask for phone numbers! \} \\
    & & Candidate Utterances:  \\
    & & \{ 1. Fang Sijin: First, change your clothes, then head to this address. A decoration company will be coming \\ 
    & &  over shortly. You’ll need to supervise their work and see how you can help. \}\\
    & & \{ 2. Zhu Shanshan: Wait, am I really responsible for this? I don’t know anything about decoration. \} \\
    & & Target Utterance: \\
    & & \{ 2. Zhu Shanshan: Wait, am I really responsible for this? I don’t know anything about decoration. \} \\
    & & Target emotion labels and cause index(es): \\
    & & [Suprise, 1] \\
    & & Input Conversation : \\
    & & ...... \\
    & & Candidate Utterances: \\
    & & ...... \\
    & & Target Utterance: \\
    & & ...... \\
    & & Target emotion labels and cause index(es): \\
    & & ...... \\
    \cline{2-3}
     Output & output example & [Happy, 1, 2] \\
  \bottomrule
  \end{tabular}
\end{table*}

\subsection{Case Study}

% 为了更加直观展示M3HG模型的优越性和局限性，我们通过两个MECAD中的对话示例展示M3HG和另外两个具有代表性的模型的预测结果。
To demonstrate the superiority and limitations of M\textsuperscript{3}HG, we present a case study that compares the prediction results of M\textsuperscript{3}HG with those of two other representative models (i.e. SHARK, GPT-4o (5-shots)), using two sample conversations from the MECAD dataset. 
As shown in Figure~\ref{fig:case study}, the first test sample demonstrates that M\textsuperscript{3}HG outperforms the other models in prediction accuracy, while GPT-4o exhibits the poorest performance.
This can be attributed to M\textsuperscript{3}HG’s use of a multimodal heterogeneous graph and a specially designed conversation super-node, which effectively captures global contextual information. These features enable M\textsuperscript{3}HG to more accurately handle scenarios where the cause utterance appears after the emotion utterance.

In the second sample, M\textsuperscript{3}HG (T+A+V) is less effective than M\textsuperscript{3}HG (T) in predicting Utterance 6 as “Sad” and Utterance 2 as “Neutral”. 
This is because the combination of text and context in Utterance 2 conveys the speaker’s worried and fearful mood, while the video and audio signals suggest a calmer demeanor.
This discrepancy likely caused M\textsuperscript{3}HG (T+A+V) to mispredict the emotions in this case. 
Nevertheless, M\textsuperscript{3}HG still outperforms all other baseline models, demonstrating its robustness and superior predictive capability even under challenging conditions.

\section{Prompt Design for ChatGPT}
\label{gpt}

We use the GPT-4o model of OpenAI public API (version up to May 13, 2024) and design a prompt elaborately to test the performance on the MECTEC task. The prompt (i.e., the input of ChatGPT) includes three parts:
\begin{itemize}
    \item \textbf{Instruction}. We use instructions to guide the ChatGPT on what it needs to do. Our instruction is as follows:

    \textit{You are an expert in sentiment analysis and identification of emotional causes. I will give you a conversation between two or more speakers. You need to extract the utter-cause-emotion triplet of the given utterance.}

    Meanwhile, we provide a detailed description of the output formats required for ChatGPT, as illustrated in Table~\ref{tab:prompts}.

    \item \textbf{Demonstrations} We achieve the few-shot in-context learning of ChatGPT by adding demonstrations. 
    We use the 5-shot in-context learning due to the limitations of the input length. 
    Each demonstration includes a conversation as input and a target utterance as the target for prediction.

\end{itemize}
Except for the aforementioned two parts, we also need to describe the conversations to be predicted and the corresponding target utterance.
An example is shown in Table~\ref{tab:prompts}.
\end{document}
